% ============================================================
% Carbon-Aware LLM Training via Hysteresis Policy Optimization
% ICREC 2026 -- Springer LNCS
% ============================================================

\documentclass[runningheads]{llncs}

\usepackage{graphicx}
\graphicspath{{./assets/}}
\usepackage{amsmath,amssymb}
\usepackage{booktabs}
\usepackage{color}
\usepackage[utf8]{inputenc}

% If you use hyperref
% \usepackage[colorlinks=true,
%   linkcolor=blue,
%   citecolor=blue,
% urlcolor=blue]{hyperref}

\usepackage{hyperref}
\usepackage{tabularx}

% ------------------------------------------------------------------
\begin{document}

\title{Carbon-Aware LLM Training via\\Hysteresis Policy Optimization}

\author{Simon Wittmann \and Julian Kraus \and Bastian
Wiesner \and Saeid Jahandar \and Reinhard German}

\authorrunning{S. Wittmann et al.}

\institute{
  Friedrich-Alexander University Erlangen-Nürnberg, Germany\\
  \email{\{sim.wittmann,boolean.kraus,bastian.wiesner,saeid.jahandar,
  reinhard.german\}@fau.de}
}

\maketitle

% ============================================================
% Abstract
% ============================================================
\begin{abstract}
%The abstract should briefly summarize the contents of the paper in 150--250 words.

  Large language model training carries a critical environmental
  cost: single runs consume millions of GPU-hours and emit hundreds of
  tonnes of CO$_2$. Temporal workload shifting can reduce these
  emissions, but existing approaches rely on ad-hoc thresholds
  without systematic policy exploration. We present a framework that
  searches the three-dimensional space of pause threshold, resume
  threshold, and start date via adaptive grid-search optimization.
  Evaluating two frontier models (DeepSeek~V3, Kimi~K2) across five
  grid regions on 2025 Electricity Maps data at 5-minute resolution,
  we find that the optimal hysteresis margin never exceeds
  16~gCO$_2$eq/kWh: wide hysteresis degrades the savings-to-overhead
  ratio, making pausing effectively a binary decision. In
  high-variability grids (Germany, Italy), savings reach 43.4\% and
  32.7\% at 174\% and 199\% time overhead. Low-variance grids (United
  States, China) yield under 10\% savings. Our results provide
  actionable, region-specific policies for carbon-aware LLM training.

  \keywords{Carbon-aware computing \and LLM training \and Hysteresis policy \and
    Grid carbon intensity \and Temporal workload shifting \and
  Simulation \and Sustainable AI.}
\end{abstract}

% ============================================================
% 1. Introduction
% ============================================================
\section{Introduction}

The energy consumption of large language model (LLM) training has
grown dramatically. Training a single frontier model requires
millions of GPU-hours: DeepSeek~V3 (2.79M~GPU-hours)
~\cite{deepseek-ai_deepseek-v3_2024} and Kimi~K2 (15.5T tokens)
~\cite{kimi_team_kimi_2025}. At a per-GPU power draw of 700~W
for an NVIDIA H800 GPU during training \cite{h800_specs} and a power
usage effectiveness (PUE) of 1.27 \cite{jegham_hungry_2025}, a single
run can emit hundreds of tonnes of CO$_2$ depending on the local grid
mix. As LLM deployment scales, managing this environmental cost
becomes essential for sustainable AI infrastructure.

Carbon-aware temporal workload shifting pauses training when grid
carbon intensity exceeds a threshold and resumes when it falls below
a second threshold, allowing a training run to consume cleaner energy
without hardware modification. However, the threshold pair
$(\theta_p, \theta_r)$
and start date $t_s$ jointly determine the savings-overhead trade-off.

Temporal workload shifting can reduce operational carbon emissions of
cloud batch jobs by 3--34\%~\cite{wiesner_lets_2021}, and
geo-distributed training during renewable curtailment windows achieves
88--95\% savings across multiple sites~\cite{wiesner_distributed_2026}.
At the datacenter scale, Carbon Explorer balances operational and
embodied carbon across hardware procurement and
deployment~\cite{acun_carbon_2023}. The emissions accounting
methodology of Patterson~et~al.~\cite{patterson_carbon_nodate} and
marginal intensity-based accounting of
Dodge~et~al.~\cite{dodge_measuring_2022} provide essential
foundations. However, none of these works provides a systematic method for
optimizing the hysteresis policy parameters of a single-site LLM
training run. Ad-hoc thresholds predominate, and the value of the
hysteresis margin (the gap between pause and resume thresholds)
remains unexplored.

  \begin{table}[htb]
    \centering
    \caption{Comparison of prior workload shifting approaches.}
    \small
    \begin{tabular}{lcc}
      \toprule
      Study & Approach & Savings \\
      \midrule
      Wiesner et al.~\cite{wiesner_lets_2021}
      & Temporal shifting, generic cloud jobs
      & 3--34\% \\
      Wiesner et al.~\cite{wiesner_distributed_2026}
      & Geo-distributed, curtailment windows
      & 88--95\% \\
      This work
      & Temporal hysteresis, frontier LLMs
      & 2--43\% \\
      \bottomrule
    \end{tabular}
  \end{table}

Our contributions to address this gap are:
  \begin{enumerate}
    \item A simulation framework for carbon-aware LLM training with
      hysteresis-based pausing, validated through structural monotonicity
      checks and cross-model scaling consistency.
    \item An adaptive grid-search optimizer that efficiently explores the
      $(\theta_p, \theta_r, t_s)$ policy space, converging within at most six
      iterations.
    \item A comprehensive empirical study covering two state-of-the-art models
      across five grid regions, establishing systematic
      Pareto frontiers, to the best of our knowledge, for
      carbon-aware LLM training.
    \item An actionable design rule: the optimal hysteresis policy uses a
      narrow margin ($\theta_p - \theta_r \leq 16$~gCO$_2$eq/kWh). Wide
      hysteresis is never Pareto-optimal.
  \end{enumerate}

  % ============================================================
  % 2. Problem Formulation
  % ============================================================
  \section{Problem Formulation}

  \subsection{System Model}

  We consider a single training job defined by: model parameter count $N$,
  token count of the dataset $T$, GPU count $G$, per-GPU training power $P_t$,
  per-GPU idle power $P_i$, facility PUE $\eta$, checkpoint save duration
  $\tau_c$, and checkpoint resume duration $\tau_r$. The job runs on a
  cluster in a single region whose grid carbon intensity $carbon(t)$ is a time
  series at a resolution of 5-minutes.

  \subsection{Hysteresis Policy}

  A hysteresis policy is defined by two thresholds and a start date:
  \begin{itemize}
    \item $\theta_p$: the pause threshold. When training and $carbon(t) > \theta_p$,
      training pauses, a checkpoint is saved, and GPUs enter idle state.
    \item $\theta_r$: the resume threshold. When paused and $carbon(t) \le \theta_r$,
      training resumes from the last checkpoint.
    \item $t_s$: the date and time at which training begins.
  \end{itemize}
  The policy is feasible only if $\theta_r \le \theta_p$; the
  difference $\Delta_\theta = \theta_p - \theta_r$ is the
  hysteresis margin. A zero margin ($\theta_r = \theta_p$) means
  the job switches on and off at a single threshold; a positive margin introduces
  a dead zone where the job remains in its current state.

  \subsection{Objective}

  Savings and overhead are computed relative to a never-pause baseline:
  \begin{equation}
    S = \frac{GHG_{b} - GHG_{p}}{GHG_{b}} \times 100 \qquad
    O = \frac{T_p - T_b}{T_b} \times 100
  \end{equation}
  where $GHG_{b}, T_b$ are the baseline (greenhouse gas) emissions and time, and $GHG_{p}, T_p$ are the policy
  emissions and time. We use a composite score that rewards carbon savings while
  enforcing an overhead budget $B$:
\[
\text{score} = \frac{
  \alpha \cdot \frac{\text{S}}{100}
  + (1-\alpha) \cdot \left(1 - \frac{\text{O}}{B}\right)
}{2}
\]

  % ============================================================
  % 3. Simulation Framework
  % ============================================================
  \section{Simulation Framework}

  \subsection{Discrete-Event Simulator}

  We implement a discrete-event simulator that replays a training run under a
  given policy. At each 5-minute time step $\Delta t$:
  \begin{itemize}
    \item If training is active and $carbon(t) > \theta_p$, the job enters a
      checkpoint-pause cycle: it spends $\tau_c$ saving state at full power,
      then remains idle at $P_i$.
    \item If paused and $carbon(t) < \theta_r$, the job resumes: it spends
      $\tau_r$ loading state (negligible with asynchronous checkpointing), then
      resumes training.
    \item Otherwise, the job continues in its current state.
  \end{itemize}
  The energy consumed per timestep is $E(t) = G \cdot P \cdot \eta \cdot \Delta t$ where
  $P$ is $P_t$ when training and $P_i$ when paused. The emissions are $GHG(t) = E(t) \cdot
  carbon(t)$. Total training tokens are tracked to ensure completion; pause
  episodes extend the duration of wall-clock time proportionally.

  \subsection{Carbon Intensity Data}

  \begin{table}
    \centering
    \caption{Grid carbon intensity characteristics (2025). Regions ordered
    by coefficient of variation.}
    \label{tab:grid}
    \small
    \begin{tabular}{lrrr}
      \toprule
      Region & Mean & Std.\ dev. & CV \\
      & (gCO$_2$eq/kWh) & (gCO$_2$eq/kWh) & \\
      \midrule
      DE (Germany) & 380 & 141.98 & 0.37 \\
      SE (Sweden)  &  23 &   8.19 & 0.35 \\
      IT (Italy)   & 309 &  87.02 & 0.28 \\
      US (United States) & 412 & 44.68 & 0.11 \\
      CN (China)   & 519 &  55.68 & 0.11 \\
      \bottomrule
    \end{tabular}
  \end{table}

  Carbon intensity data comes from the Electricity Maps API at
  5-minute resolution for five regions in 2025. Table~\ref{tab:grid}
 summarizes the mean intensity, standard deviation, and
  coefficient of variation (CV) of each region.
  Grid variability determines savings potential: Germany
  and Italy offer the greatest temporal arbitrage
  opportunity, while the United States and China offer few
  clean intervals. The time series exhibits strong temporal
  autocorrelation (lag-1 $= 0.999$, one-day lag $= 0.694$ for
  Germany), which makes a chronological trace replay essential.

  % ============================================================
  % 4. Optimization Method
  % ============================================================
  \section{Adaptive Grid-Search Optimization} \label{sec:opti-algo}

  Exhaustively enumerating the 3D policy space $(\theta_p, \theta_r, t_s)$ is
  computationally prohibitive: a naive grid at resolution $N$ per axis costs
  $O(N^2 \cdot D)$ where $D$ is the number of start-date samples. With
  $N=100$ and $D=52$, brute-force requires approximately 260,000 simulations.
  We instead use an adaptive grid-search algorithm inspired by trust-region
  optimization:
  \begin{enumerate}
    \item {\bf Initialization.} A coarse grid over full bounds:
      % obacht, this is correct!
      $\theta_p \in [0, \theta_p^{\max}]$, $\theta_r \in [0, \theta_p^{\max}]$ with $\theta_r \le \theta_p$, and
      $t_s$ at $D$ uniform intervals. The maximum $\theta_p^{\max}$ is
      region-dependent: 100~gCO$_2$eq/kWh for Sweden, 800 for all others.
    \item {\bf Evaluation and selection.} Every grid point is simulated. Points
    that exceed the overhead budget $B$ are discarded.
    \item {\bf Zoom refinement.} Bounds contract around the best point by
    shrinking factor $s = 0.45$ per axis. The grid resolution $N = 10$ points per
      axis per iteration remains fixed. If the contracted interval on any axis
      falls below $3$~gCO$_2$eq/kWh, the search on that axis terminates.
    \item {\bf Iteration.} Steps 2--3 repeat for a maximum of $K = 10$
      iterations.
  \end{enumerate}

  The algorithm guarantees a feasible point (the initial grid always contains
    the no-pause policy $\theta_p = \theta_p^{\max}, \theta_r =
  0$, yielding $0\%$ overhead). The total complexity is $O(K \cdot N^2 \cdot D)$;
  with our settings ($K=10, N=10, D=7$ for start-date experiments, $D=1$ for
  fixed-date runs), a full optimization evaluates approximately 3,500
  candidate policies (550 for fixed-start runs), completing in under one
  minute on a single CPU core. The optimizer converges within 5--6
  iterations in all (model, region) combinations; subsequent
  iterations yield no measurable improvement
  ($\Delta\text{score}\leq2\times10^{-5}$).

  % ============================================================
  % 5. Experimental Setup
  % ============================================================
  \iffalse
  \section{Experimental Setup}

  \subsection{Models}

  We evaluated DeepSeek~V3 (671B total parameters, 37B active with MoE,
  trained on 14.8T tokens using 2048~H800~GPUs, 2.79M~GPU-hours) and
  Kimi~K2 (1T parameters, 32B active with MoE, trained on 15.5T tokens,
  minimum 256~GPU model-parallel group). The GPU-hours for Kimi were
  estimated based on the GPU-hours per trillion tokens given by Deepseek.
  These represent the current frontier of LLM pretraining at
  substantially different scales.

  \subsection{Hardware Constants}

  GPU training power: 700~W (NVIDIA H800 TDP~\cite{h800_specs}). GPU idle
  power: 60~W (estimated). PUE: 1.27 (DeepSeek DGX H800 infrastructure
  data~\cite{jegham_hungry_2025}). Checkpoint pause: 148.8~s (PyTorch
  Distributed synchronous checkpoint, 7B model~\cite{pytorch_checkpoint}).
  Checkpoint resume: 0~s (asynchronous checkpoint I/O).

  \subsection{Regions and Data}

  Five regions shown in Table~\ref{tab:grid}. We use 2025 marginal carbon
  intensity at 5-minute resolution from Electricity Maps. Our framework
  generalizes to any year of available intensity data.

  \subsection{Optimizer Settings}

  Resolution $N=10$, iterations $K=10$, budget $B=200\%$, maximum
  pause threshold $\theta_p^{\max}$, shrink factor
  $s=0.45$, minimum step $3$~gCO$_2$eq/kWh, $\alpha=1.0$ (pure savings
  objective within budget). Start-date resolution: 1 sample (Jan~1) for
  fixed-date runs, 7 samples per year for start-date optimization.
\fi

  % ============================================================
  % 6. Results
  % ============================================================
  \section{Results}

  \subsection{Experimental Setup}

\begin{table}[htbp]
  \centering
  \caption{Experimental Setup used in our simulation runs.}
  \label{tab:exp_setup}
  \small
  \begin{tabularx}{\columnwidth}{|l|X|}
    \toprule
    \textbf{Parameter / Setting} & \textbf{Specification / Value} \\
    \midrule
    \multicolumn{2}{|l|}{\textbf{Models}} \\
    DeepSeek V3 & 671B total parameters (37B active MoE), 14.8T tokens, 2048\,H800 GPUs, 2.79M GPU-hours \\
    Kimi K2 & 1T total parameters (32B active MoE), 15.5T tokens, $\ge~256$ GPU model-parallel group\textsuperscript{*} \\
    \midrule
    \multicolumn{2}{|l|}{\textbf{Hardware Constants}} \\
    GPU Power (Active / Idle)& 700\,W (H800 TDP)~\cite{h800_specs} \quad / 60\,W (est.) \\
    PUE & 1.27 (DeepSeek DGX H800 infrastructure)~\cite{jegham_hungry_2025} \\
    Checkpoint Overhead- & \\
    Time (Pause / Resume) & 148.8\,s (PyTorch Dist. 7B)~\cite{pytorch_checkpoint} / 0\,s (async I/O)\\
    \midrule
    \multicolumn{2}{|l|}{\textbf{Additional Optimizer Settings (not described in section \ref{sec:opti-algo})}} \\
    Budget \& Weight & Budget $B=200\%$, $\alpha=1.0$ (pure savings objective) \\
    \bottomrule
  \end{tabularx}
  \vspace{1mm}
  
  \raggedright \scriptsize \textsuperscript{*}GPU-hours estimated based on DeepSeek GPU-h/trillion tokens.
\end{table}

  \subsection{Pareto Frontier Analysis}

  Figure~\ref{fig:pareto} presents the combined Pareto frontiers.
  Table~\ref{tab:best} reports the best-score configurations. For
  DeepSeek~V3, single-year savings range from 43.4\% (Germany, score 0.717)
  to 6.2\% (China, score 0.531).
  
  In high-variability grids, the frontier reflects substantial carbon-intensive intervals 
  worth avoiding. Sweden plateaus earlier: its consistently low baseline carbon intensity yields fewer high-emission windows to mitigate. The United States and China show limited 
  temporal variation: pausing rarely finds a significantly cleaner interval 
  than the one it left. 

  \begin{figure}[h]
    \centering
    \includegraphics[width=0.7\textwidth]{savings_vs_overhead_all}
    \caption{
      Combined Pareto frontiers for DeepSeek~V3 and
      Kimi~K2 across five grid regions.
      Each point is a feasible policy $(\theta_p, \theta_r, t_s)$
      evaluated on 2025 5-minute intensity data.
      Best-score points within 200\% overhead are stars.
    }
    \label{fig:pareto}
  \end{figure}


\begin{table}[htbp]
  \centering
  \caption{Best policy configurations per (model, region) with optimized
    start dates. The hysteresis margin $\Delta_\theta = \theta_p -
    \theta_r$ is
  $\leq 16$~gCO$_2$eq/kWh across all combinations.}
  \label{tab:best}
  \small
  \begin{tabular}{l | rrrr | rrrr}
    \toprule
    & \multicolumn{4}{c|}{DeepSeek~V3} & \multicolumn{4}{c}{Kimi~K2} \\
    \cmidrule(lr){2-5} \cmidrule(lr){6-9}
    Region & $\theta_p$ & $\theta_r$ & Savings & Overhead & $\theta_p$ & $\theta_r$ & Savings & Overhead \\
    \midrule
    DE & 272 & 268 & 43.4\% & 174.3\% & 267 & 262 & 27.5\% & 199.4\% \\
    IT & 246 & 231 & 32.7\% & 194.7\% & 283 & 271 & 16.6\% & 129.1\% \\
    SE &  19 &  18 & 23.2\% & 106.3\% &  22 &  21 & 12.8\% & 115.4\% \\
    US & 396 & 386 &  9.5\% & 115.7\% & 453 & 447 &  3.1\% &  22.2\% \\
    CN & 528 & 512 &  6.2\% &  72.5\% & 537 & 522 &  2.4\% &  40.1\% \\
    \bottomrule
  \end{tabular}
\end{table}

  For Kimi~K2, the same regional ordering holds at lower absolute savings
  (Germany 27.5\%, score 0.637; China 2.4\%, score 0.512). Larger GPU counts
  amplify the benefit of avoiding carbon intensive intervals, but the scaling is
  sub-linear: DeepSeek~V3's up to 8$\times$ GPU advantage yields only
  1.5--2.0$\times$
  the relative savings, because checkpoint overhead scales with GPU count.

  The ranking of regions by savings potential (DE $>$ IT $>$ SE $>$ US $>$ CN) follows
  the ranking by coefficient of variation; mean intensity plays a secondary
  role. Sweden's high coefficient of variation is offset by its very low mean (23~gCO$_2$eq/kWh),
  limiting the absolute emissions available to avoid.

 

  \subsection{Threshold Space and Hysteresis Margin}

  \begin{figure}
    \centering
    \includegraphics[width=0.55\textwidth]{threshold_space_DeepSeek_V3_DE_0101}
    \caption{
      Pause vs.\ resume threshold for DeepSeek~V3 training in Germany starting on July 1st.
      Color encodes composite score. The diagonal $\theta_r = \theta_p$
      marks zero-hysteresis policies.
    }
    \label{fig:threshold}
  \end{figure}

  The threshold space analysis (example in Fig.~\ref{fig:threshold})
  reveals the central
  finding: optimal hysteresis margins cluster near zero. Across all 10
  (model, region) combinations, the best-score points satisfy $\Delta_\theta
  \leq 16$~gCO$_2$eq/kWh, and in most cases $\Delta_\theta \leq
  10$~gCO$_2$eq/kWh (Table~\ref{tab:best}). Zero-hysteresis policies
  ($\theta_r = \theta_p$) achieve scores within 1--3\% of the optimum,
  confirming that a non-zero margin provides marginal practical benefit.


  \subsection{Start Date Sensitivity}

 Although the start date exerts only a secondary influence in some zones, its impact is substantial in others. Across all (model, region) combinations, winter months (November--February) are
  disproportionately represented among optimal start dates, reflecting lower
  grid carbon intensity during the heating season. The effect ranges from
  approximately 20 percentage points of savings (Italy) to under 2 percentage
  points (China), confirming that the hysteresis threshold choice mostly, but not always, dominates
  start-date scheduling.

  % ============================================================
  % 7. Discussion
  % ============================================================
  \section{Discussion and Conclusion}



  \subsection{Limitations}
  %\paragraph{Limitations}

  Our simulation assumes perfect foresight: policies are evaluated on
  historical data rather than forecasted future intensity. A real deployment
  would require forecasting, introducing prediction error. We use a fixed
  checkpoint overhead, real overhead depends on model size and
  storage bandwidth. The idle power of the GPU is an estimate; the actual values vary.
  Finally, this work considers single-site training; geo-distributed training
  could compound savings by adding spatial flexibility.

  % ============================================================
  % 8. Conclusion
  % ============================================================
  \subsection{Conclusion}
  %\paragraph{Conclusion}

  We have presented a simulation framework for carbon-aware
  LLM training via hysteresis-based pausing. Across two models and five
  grid regions, it establishes systematic Pareto
  frontiers, to the best of our knowledge, for this problem.

  The central finding is that hysteresis provides diminishing returns: the
  optimal policy uses a narrow or zero margin ($\Delta_\theta \leq
  16$~gCO$_2$eq/kWh). Wide hysteresis margins waste
  time and checkpoint overhead without proportional savings. Carbon-aware
  pausing can reduce operational emissions by up to 43\% in high-variability
  grids at the cost of up to 200\% wall-clock overhead.

  In practice: in high-variability grids, set $\theta_p$ at the
  20--40th percentile of regional intensity with near-zero hysteresis
  for 30--43\% savings at up to 200\% overhead. Low-variance grids
  yield under 10\% savings; geo-distributed training or embodied
  carbon trade-offs are more promising there.

  Three directions are particularly promising for future work: (1) extending
  the framework to use forecasted rather than historical carbon intensity for
  real-time deployment; (2) combining temporal shifting with geo-distributed
  training; and (3) incorporating embodied carbon awareness to ensure
  operational savings are not outweighed by manufacturing emissions.
  Beyond individual emissions reduction, carbon-aware scheduling also
  provides demand-side flexibility that supports grid integration of
  renewable energy.

  All simulation code and data are available on
  Github\footnote{https://404simon.github.io/TheGreenEpoch},
  where our interactive browser-based simulator and optimizer can be tried.

  % ============================================================
  % References
  % ============================================================
  \bibliographystyle{splncs04}
  \bibliography{references}

  \end{document}
